\section{How much can one divide out of a period?}

% BEGIN I89M1
\paragraph{I89M1. Both word modules are full lattices.}
Fix any guarded family member $D=\mathcal F.D(x,i)$, $x>\mathcal F.x_1$,
a word $u$, its level prime $T$, and the original coefficient field
$K$ with ring of integers $O$. Let $V$ be the rational word space,
$M$ its ramified lattice, and $N=L.\mathrm{NuLat}(u)$ its unshifted
module. We prove that $N$ is finitely generated and that
$cM\subset N$ for some nonzero $c\in O$.
The shifted module in the divisibility problem is therefore full as well.

Write $\Lambda$ for the given full lattice at stage zero. The injection
$\nu:V\to V_{\rm stage}$ identifies $M$ with the intersection of
$\Lambda$ and a rational subspace. Thus $M$ is finite over $O$ and
spans $V$, by clearing denominators of a rational basis. For every
$w\nmid T$, the defining lattice condition puts $N_w\subset M_w$.
Here and below the subscript denotes localization inside the rational
space. Choose a prime $p\ne T$ where the unshifted lattice is required
to equal $M$ locally. In the host, $2$ and the distinguished odd prime
$q$ both lie outside its two inert or signal classes, so one can take
$2$ unless $T=2$, and $q$ otherwise. At any $w\mid p$ the family
condition gives $N_w=M_w$. For each of finitely many generators of
$M$ clear its local denominator at this $w$. The product of the
denominators gives $cM\subset N$.

It remains to bound the possible denominators of $N$ at $w\mid T$.
The family defines $N_w$ as the inverse image, under $\nu$, of a
module obtained by applying a specified complex Hecke operator to
the localized ramified input. More explicitly let $j$ be the
injective complex realization of $V_{\rm stage}$, and let $A_w$
be the intersection of $\Lambda_w$ with the ramified rational
subspace. With the original scalar power of $T$ included in the
operator $H_w$, this definition reads
\[
       z\in N_w\quad\Longleftrightarrow\quad
       j(\nu z)\in H_w(j(A_w)).
\]
The DVR $O_w$ is Noetherian. Hence $A_w\subset\Lambda_w$ is finite
over $O_w$. The complex linearity of $H_w$ and the prescribed
coefficient embedding make $H_w(j(A_w))$ finite over $O_w$ too, since the images of any generators suffice. Intersect this finite module
with $j(\nu V)$ and pull it back through the injection $j\nu$.
The intersection is a submodule over a Noetherian ring. This proves
that $N_w$ is finitely generated, even before establishing any
integrality estimate for the Hecke operator.

Since $M_w$ spans $V$, a power of $T$ carries these finitely many
local generators into $M_w$. Take the largest exponent $b\geq0$
at the finitely many $w\mid T$. At other places $N_w\subset M_w$
already, so $T^bN\subset M$ by localization at all maximal ideals.
Now $N$ is a submodule of the finite $O$-module $T^{-b}M$, hence
finite. Together with $cM\subset N$ this proves the assertion.
Finally the native shifted module is $I_uN$, where $I_u$ is the
specified finite product of nonzero fractional prime-ideal powers.
It is finitely generated and spans $V$. This gives the original
$\mathrm{NuCommens}$ and $\mathrm{NuFull}$ at every guarded member,
hence both eventual family conclusions with threshold
$\mathcal F.x_1+1$. The exponent $b$ and element $c$ are allowed to
depend on the member and word. The local ring and finite-module facts
used here are the Dedekind criterion and Noetherianity
(Stacks Project, Algebra, Lemma 10.120.17, Tag 034X).
% END I89M1

% BEGIN I89M2
\paragraph{I89M2. The valuation to test is the content of a pairing.}
On the selected class of R88L1--3, let $R$ be the sharp row lattice
and let $A_t$ be the module generated by the rational rows realized
at tuple $t$, as in B87H1. For every attained $t$ it is finite and
spans the rational row space $W$. In particular the entire family of
Hermitian lattices satisfies the original $\mathrm{AllGood}$ condition because
I89M1 gives fullness of the word lattices, R88L2 gives fullness of
the sharp row lattice, and the member already specifies invertible
metric frames at every infinite place. Put $t_w=\mu_{\rm At}(w)$.
The family defines $R=\bigcap_w(A_{t_w})_w$ inside $W$.
In fact its localization at each $w$ is the entire
$(A_{t_w})_w$. To see this, take $y\in(A_{t_w})_w$.
Only finitely many other places $v$ can have
$y\notin(A_{t_v})_v$, since the base lattice is full, and the tuple differs
from the base only on the finite support. At each of those
places $v$ some power of its prime ideal carries $y$ into the
local lattice. Their product is coprime to $w$, so it contains
an element $s$ outside $w$. Then $sy\in R$. This proves the
claimed equality of localizations.

Let $B_u:V\times W\to K$ be the actual native rational pairing.
For finite generators $n_a$ of the unshifted module $N$ and $r_{tb}$
of $A_t$, form the fractional ideal
\[
     C_{u,t}=\sum_{a,b} O\,B_u(n_a,r_{tb}).
\]
Use valuation $+\infty$ for the zero ideal. This ideal is
independent of generators by bilinearity. Put
$k_{u,w}=\lfloor e_w\,\mathrm{sg}(u,w)\rfloor$ on the original
support and $k_{u,w}=0$ off it. The values on the actual two localized
lattices generate exactly
\[
   B_u\big((I_uN)_w,R_w\big)\,O_w
        =\mathfrak p_w^{-k_{u,w}} C_{u,t_w}O_w ,
\]
where the left side means the $O_w$-span of all paired values.
Indeed the other factors of $I_u$ are units at $w$ and we just
computed $R_w$. A functional belongs to the integral dual if its
values on the lattice are in $O$. Testing this at all finite places
now gives the equivalence
\[
   \mathrm{OneMapIntoCC}\quad\Longleftrightarrow\quad
   v_w(C_{u,t_w})\ge k_{u,w}\quad\text{for every }u,w .
\]
In particular the exact exponent of any extra local denominator at
$w$ is $\max\{0,k_{u,w}-v_w(C_{u,t_w})\}$, interpreted as zero
for a zero pairing. An overall denominator does not establish the
original entry unless each of these exponents vanishes.
% END I89M2

% BEGIN I89M3
\paragraph{I89M3. All but finitely many places satisfy the integral entry on each member.}
For the base tuple $t_0$ let
$J_u=\{a\in O:aC_{u,t_0}\subset O\}$. This is a nonzero integral
ideal since I89M1 and B87H1 give finite generating sets. It is $O$ if
$C_{u,t_0}=0$. For the given member it suffices to test only the
primes in
\[
 L.\mathrm{supp}\ \cup\
 \{w:\mathfrak p_w\mid\prod_u J_u\}.
\]
Outside this set $k_{u,w}=0$, $t_w=t_0$, and every value on the
unshifted lattice and the base row lattice is already in $O_w$.
There are finitely many words and hence this gives simultaneous
integrality at every other place. At a listed place the minimum
valuation of the entries of the finite matrix in I89M2 gives
the exact comparison. This finite exceptional set depends on the
member. The original family theorem requires the inequalities
even at its exceptional places for every sufficiently large member.

There is a second way to prove those inequalities. In the original
coordinate construction the value of a period on a good integral
column is a local unit times the value of $B_u$ on its row image.
Those images generate the local row lattice. A lower valuation bound
$k_{u,w}$ on all such periods, for unshifted inputs, therefore gives
the same bound for their span, by the ultrametric inequality.
In applying this route one needs the unit identity on the actual
tuples and the local lower bound on the actual unshifted inputs.
The source proves the latter by separate calculations at the word
prime $T$, at the two generator primes away from $T$, and at the
remaining primes. At $T$ it uses the ordinary and twisted
level-prime column estimates. At a generator it uses a bound on the
used tuples together with the prescribed gain inequality. Elsewhere it
compares the integral ramified input and the served local model in the
chain of column maps and subtracts the stabilizer denominator. The
selected uniform comparison still requires these estimates or an
independent evaluation of the content ideals. The finite reduction
and the fullness assertion above use neither an archimedean bound
nor equality with an analytic Bessel period.
% END I89M3
